\section{코드 블록 구현 방식 --- 웹앱 제작 시 요구사항}

본 장은 Google Blockly 기반 블록 코딩이 어떻게 정의되고, 어떻게 실행 가능한
명령으로 변환\textbf{·}전송되는지를 \code{blocklyconfig.js}와
\code{commandexecutor.js} 중심으로 분석한다.

\subsection{블록 정의 (\code{blocklyconfig.js})}

커스텀 블록은 Blockly의 JSON 스키마로 정의된다. 각 블록은 고유 \code{type},
시각적 레이블 \code{message0}, 입력/필드 정의 \code{args0}, 카테고리 색상
\code{colour}, 한국어 도움말 \code{tooltip}을 가진다.

\captionof{listing}{블록 정의 예시(개념)}
\begin{minted}{json}
{
  "type": "timed_forward",
  "message0": "🚗 서보 전진 %1 초",
  "args0": [{ "type": "input_value", "name": "SECONDS", "check": "Number" }],
  "previousStatement": null, "nextStatement": null,
  "colour": "#FF8C00",
  "tooltip": "지정한 시간 동안 앞으로 이동합니다"
}
\end{minted}

\subsection{블록 카테고리(툴박스)}

툴박스는 \code{main.html}의 XML로 정의되며, 이모지 라벨을 가진 카테고리들로
구성된다.

\begin{longtable}{L{3.4cm} L{2.6cm} L{8.0cm}}
\toprule
\rowcolor{tablehead}\textbf{카테고리} & \textbf{색상} & \textbf{대표 블록} \\
\midrule
\endhead
\emj{🚗} 서보 모터 & \#FF8C00 & \code{timed\_forward/backward/left/right}, \code{move\_forward}, \code{turn\_left/right}, \code{stop\_moving} \\
\emj{⚡} DC 모터 & \#FFCC00 & \code{main\_motor\_forward\_timed}, \code{main\_motor\_forward}, \code{main\_motor\_stop} \\
\emj{💡} LED & \#FF5555 & \code{led\_on}, \code{led\_off}, \code{led\_off\_all}, \code{set\_lamp} \\
\emj{🖥} 디스플레이 & \#9966FF & \code{send\_message}, \code{send\_message\_xy}, \code{display\_icon}, \code{clear\_display}, \code{clear\_rect} \\
\emj{🔊} 소리 & \#00CCFF & \code{buzzer\_on}, \code{buzzer\_note}(3옥타브 음계) \\
\emj{🔫} 발사 & \#FF4500 & \code{gun\_fire} \\
\emj{📡} 센서 & \#5C81A6 & \code{pico\_check\_device}, \code{check\_distance}, \code{check\_magnetic} \\
\emj{⏱} 시간 & \#5CA65C & \code{time\_sleep} \\
\emj{🔁} 제어 & (내장) & \code{controls\_repeat\_ext}, \code{controls\_whileUntil}, \code{controls\_if} \\
\emj{🚀} 한꺼번에 & \#8E44AD & \code{batch\_block} \\
\emj{📦} 변수 / \emj{🔢} 수학\textbf{·}논리 / \emj{🔧} 함수 & (내장) & Blockly 표준 블록 \\
\bottomrule
\end{longtable}

\subsection{블록 → 명령 변환 (\code{commandexecutor.js})}

\code{CommandExecutor}는 블록 트리를 순회하며 각 블록을 프로토콜 명령 문자열로
변환하고 전송한다.

\subsubsection{명령 생성: \code{generateCommand(block)}}

\begin{longtable}{L{4.4cm} L{10.0cm}}
\toprule
\rowcolor{tablehead}\textbf{블록 타입} & \textbf{생성 명령} \\
\midrule
\endhead
\code{timed\_forward} & \code{SERVO\_tFORWARD,\$\{초\}} \\
\code{move\_forward} & \code{SERVO\_FORWARD} \\
\code{main\_motor\_forward\_timed} & \code{DC\_tFORWARD,\$\{초\}} \\
\code{led\_on} & \code{LED\_ON,\$\{번호\},\$\{밝기\}} \\
\code{led\_off\_all} & \code{LED\_OFF,ALL} \\
\code{set\_lamp} & \code{[\$\{v0\} \$\{v1\} ... \$\{v5\}]} \\
\code{send\_message} & \code{MSG,\$\{로마자변환텍스트\}} \\
\code{display\_icon} & \code{ICON,\$\{이름\},\$\{x\},\$\{y\}} \\
\code{buzzer\_on} / \code{buzzer\_note} & \code{BUZZER\_ON,\$\{주파수\},\$\{지속\}} \\
\code{gun\_fire} & \code{GUN\_FIRE} \\
\code{check\_distance} & \code{DISTANCE} ($\to$ 변수 \code{\_last\_distance}) \\
\code{check\_magnetic} & \code{MAGNET} ($\to$ 변수 \code{\_last\_magnetic}) \\
\code{pico\_check\_device} & \code{PING} \\
\code{time\_sleep} & \code{SLEEP,\$\{초\}} \\
\bottomrule
\end{longtable}

OLED는 ASCII 폰트만 지원하므로, \code{send\_message}의 한글 텍스트는
\code{romanize.js}로 로마자 변환된 뒤 전송된다(예: ``안녕'' $\to$ ``annyeong'').

\subsubsection{값 블록 평가: \code{evaluateValueBlock(block)}}

숫자\textbf{·}텍스트\textbf{·}변수\textbf{·}산술\textbf{·}비교\textbf{·}논리\textbf{·}난수\textbf{·}함수 반환
블록을 재귀적으로 평가하여 문자열 값을 산출한다(0으로 나눗셈 보호 포함).
센서 값은 \code{state.variables}의 \code{\_last\_distance},
\code{\_last\_magnetic}에 저장되어 조건식에서 참조된다.

\subsubsection{블록 처리와 제어흐름: \code{processBlock(block)}}

블록 체인을 순회하며 반복(\code{controls\_repeat\_ext}), 조건(\code{controls\_if}),
while/until(\code{controls\_whileUntil}), 변수 설정을 해석한다. 무한 루프
방지를 위해 반복 횟수와 중첩 깊이에 상한을 둔다.

블록 하나가 실행 경로로 들어가 명령으로 변환\textbf{·}전송되기까지의 파이프라인을
정리하면 그림~\ref{fig:blockpipe}와 같다. 블록은 성격에 따라 세 갈래로 분기한다.
(1) \code{batch\_block}은 평탄화 후 일괄 전송, (2) 제어흐름\textbf{·}변수\textbf{·}함수
블록은 \code{handleLogicBlock}이 재귀 처리, (3) 그 외 동작 블록은
\code{generateCommand}로 명령 문자열을 만들어 전송한다. 조건식\textbf{·}인자는
\code{evaluateValueBlock}이 재귀 평가하여 값을 공급한다.

\begin{diagram}{블록 → 명령 처리 파이프라인}{fig:blockpipe}
\node[boxw, text width=26mm] (ws) {executeWorkspace\\{\scriptsize 최상위 블록 순회}};
\node[boxw, right=10mm of ws, text width=24mm] (pb) {processBlock\\(block)};
\node[boxw, above right=6mm and 12mm of pb, text width=26mm] (batch) {\_processBatch\\{\scriptsize BATCH;... 일괄}};
\node[boxw, right=14mm of pb, text width=26mm] (logic) {handleLogicBlock\\{\scriptsize 반복·조건·변수}};
\node[boxw, below right=6mm and 12mm of pb, text width=26mm] (gen) {generateCommand\\{\scriptsize 블록→문자열}};
\node[boxw, below=7mm of gen, text width=26mm] (send) {sendCommand\\$\to$ \_dispatch};
\node[hw, below left=7mm and -6mm of send, text width=20mm] (bt) {BLE 전송};
\node[hw, below right=7mm and -6mm of send, text width=20mm] (sim) {시뮬레이터};
\node[store, above=7mm of logic, text width=22mm] (ev) {evaluateValueBlock\\{\scriptsize 값 재귀 평가}};
\draw[flow] (ws)--(pb);
\draw[flow] (pb) -- node[lbl,above]{batch} (batch);
\draw[flow] (pb) -- node[lbl,above]{logic} (logic);
\draw[flow] (pb) -- node[lbl,below]{동작} (gen);
\draw[flow] (gen)--(send);
\draw[flow] (send) -- node[lbl,left]{실기} (bt);
\draw[flow] (send) -- node[lbl,right]{시뮬} (sim);
\draw[dep] (ev) -- node[lbl,left]{값} (logic);
\draw[dep] (ev.south west) to[bend right=10] (gen.north east);
\draw[flow] (logic.south) to[bend left=10] node[lbl,right]{재귀} (pb.east);
\end{diagram}

\subsubsection{명령 디스패치와 Fire-and-forget}

\code{sendCommand()}는 명령 헤드가 \code{FIRE\_AND\_FORGET\_HEADS} 집합에
속하는지로 응답 대기 여부를 결정한다. 시뮬레이션 모드에서는 BLE 대신
시뮬레이터 싱크(\code{simSink})로 명령을 보낸다.

\begin{itemize}
  \item \textbf{응답 불필요}(fire-and-forget): \code{LED\_*}, \code{MSG*}, \code{ICON},
        \code{SERVO\_*}/\code{DC\_*} 연속 구동, \code{BUZZER\_ON}, \code{SYS\_SET} 등.
  \item \textbf{응답 대기}: \code{DISTANCE}, \code{MAGNET}, \code{PING}, \code{SLEEP},
        \code{tFORWARD} 계열(시간 명령), \code{BATCH}.
\end{itemize}

\subsection{한꺼번에 실행: \code{batch\_block}}

\code{batch\_block}은 내부 블록들을 평탄화하여 \code{BATCH;cmd1|cmd2|...}
단일 페이로드로 전송하고, 피코로부터 단일 ACK를 받는다. 센서\textbf{·}제어흐름\textbf{·}
변수\textbf{·}함수 블록은 배치 내부에 넣을 수 없으며(\code{BATCH\_FORBIDDEN\_TYPES}),
\code{onchange} 검증기가 위반 시 경고한다.

\subsection{웹앱 제작 시 요구사항}

\begin{infobox}[title=새 블록 추가 시 변경 지점]
새로운 블록을 추가하려면 일반적으로 세 곳을 함께 수정해야 한다.
\begin{enumerate}
  \item \code{blocklyconfig.js} --- 블록 정의(타입, 메시지, 인자, 색상).
  \item \code{main.html} --- 툴박스 XML에 블록 등록(필요 시 shadow 기본값).
  \item \code{commandexecutor.js} --- 블록 타입 $\to$ 명령 문자열 생성 및 응답 정책.
\end{enumerate}
\end{infobox}

추가로, 새 명령이 펌웨어 동작을 요구한다면 \code{Pico/process\_data.py}의
\code{CommandProcessor}에 대응 핸들러를 추가해야 한다(9장 연계 구조 참조).
OLED로 출력되는 한글 텍스트는 로마자 변환 경로를 거치도록 하고, 다중 청크가
되는 긴 명령은 BLE 청크 타이밍을 고려해야 한다.

\begin{teachbox}{블록에서 명령으로}
\teachhead{설계 의도}
\begin{itemize}
  \item 블록을 batch/제어흐름/동작 \emph{세 갈래}로 나눠 처리한다.
  \item 명령 생성 로직을 \emph{하나}로 두어, 시뮬레이터와 실기가 같은 의미로 동작한다.
\end{itemize}
\teachhead{분석할 때 핵심 질문}
\begin{itemize}
  \item 새 블록을 추가할 때 왜 \code{blocklyconfig.js}\textbf{·}툴박스\textbf{·}\code{commandexecutor.js} 세 곳을 함께 고쳐야 하는가?
  \item 시뮬과 실기가 같은 명령을 쓰면 교육적으로 무엇이 좋은가?
\end{itemize}
\teachhead{점검 요소}
\begin{itemize}
  \item 블록 1개를 정의 $\to$ 툴박스 등록 $\to$ 명령 생성까지 끝까지 추가하라(펌웨어 동작이 필요하면 핸들러도 함께).
\end{itemize}
\end{teachbox}

